\begin{lemma}
\label{lemma-completion-disconnected}
Let $(A, \mathfrak m)$ be a Noetherian local ring. The punctured spectrum
of $A^\wedge$ is disconnected if and only if the punctured spectrum of $A^h$
is disconnected.
\end{lemma}

\begin{proof}
Since the completion of $A$ and $A^h$ are the same, we may assume
that $A$ is henselian (Lemma \ref{lemma-henselization-noetherian}).

\medskip\noindent
Since $A \to A^\wedge$ is faithfully flat (see reference just given)
the map from the punctured spectrum of $A^\wedge$ to the punctured
spectrum of $A$ is surjective
(see Algebra, Lemma \ref{algebra-lemma-ff-rings}).
Hence if the punctured spectrum of $A$ is disconnected, then
the same is true for $A^\wedge$.

\medskip\noindent
Assume the punctured spectrum of $A^\wedge$ is disconnected.
This means that
$$
\Spec(A^\wedge) \setminus \{\mathfrak m^\wedge\} = Z \amalg Z'
$$
with $Z$ and $Z'$ closed. Let
$\overline{Z}, \overline{Z}' \subset \Spec(A^\wedge)$ be the closures.
Say $\overline{Z} = V(J)$, $\overline{Z}' = V(J')$ for some ideals
$J, J' \subset A^\wedge$. Then $V(J + J') = \{\mathfrak m^\wedge\}$
and $V(JJ') = \Spec(A^\wedge)$. The first equality means that
$\mathfrak m^\wedge = \sqrt{J + J'}$ which implies
$(\mathfrak m^\wedge)^e \subset J + J'$ for some $e \geq 1$.
The second equality implies every element
of $JJ'$ is nilpotent hence $(JJ')^n = 0$ for some $n \geq 1$.
Combined this means that $J^n/J^{2n}$ is annihilated by
$J^n$ and $(J')^n$ and hence by $(\mathfrak m^\wedge)^{2en}$.
Thus we may apply Lemma \ref{lemma-quotient-by-idempotent}
to see that there is a finite type $A$-algebra $C$ and an
isomorphism $A^\wedge/J^n = C^\wedge$.

\medskip\noindent
The rest of the proof is exactly the same as the second part
of the proof of Lemma \ref{lemma-one-dimensional-formal-branch};
of course that lemma is a special case of this one!
We have $C/\mathfrak m C = A/\mathfrak m$ because $A^\wedge/J^n$ has the same
property. Hence $\mathfrak m_C = \mathfrak m C$ is a maximal ideal
and $A \to C$ is unramified at $\mathfrak m_C$
(Algebra, Lemma \ref{algebra-lemma-characterize-unramified}).
After replacing $C$ by a principal localization we may
assume that $C$ is a quotient of an \'etale $A$-algebra $B$, see
Algebra, Proposition \ref{algebra-proposition-unramified-locally-standard}.
However, since the residue field extension of $A \to C_{\mathfrak m_C}$
is trivial and $A$ is henselian, we conclude that $B = A$
again after a localization.
Thus $C = A/I$ for some ideal $I \subset A$ and it follows that
$J^n = IA^\wedge$ (because completion is exact in our situation by
Algebra, Lemma \ref{algebra-lemma-completion-flat}) and $I = J^n \cap A$
(by flatness of $A \to A^\wedge$).
By symmetry $I' = (J')^n \cap A$ satisfies $(J')^n = I'A^\wedge$.
Then $\mathfrak m^e \subset I + I'$ and $II' = 0$
and we conclude that $V(I)$ and $V(I')$ are closed subschemes
which give the desired disjoint union decomposition of the
punctured spectrum of $A$.
\end{proof}